\chapter{Random matrices and zeta zeros}\label{ch:oc-gue}

This chapter states Montgomery's pair correlation for the zeros of $\zeta$, the comparison with the Gaussian unitary ensemble, and the numerical test Odlyzko carried out.

\section{Montgomery's pair correlation}

Assume the Riemann hypothesis and write the ordinates of the zeros as $\frac12+i\gamma$. Montgomery's theorem, for test functions whose Fourier transform is supported in $[-1,1]$, says that the pair correlation of the $\gamma$ matches the GUE form $1-\bigl(\frac{\sin\pi u}{\pi u}\bigr)^2$.

\section{GUE statistics}

The Gaussian unitary ensemble is the random Hermitian matrix with density proportional to $e^{-\operatorname{tr} H^2/2}$. Its eigenvalue spacings are the model for the normalised gaps between zeros.

\section{Odlyzko's computations}

Odlyzko computed millions of zeros at heights past $10^{20}$ and compared their spacing distribution with the GUE prediction. The agreement is numerical evidence, not a proof.
